% ============================================
% DOC-26 — Anotações de Oliveira — O Mecanismo (edição revista R1075)
% ============================================
\fragtitle{doc}{D}{DOC}{26}{Anotações de Oliveira --- O Mecanismo}
\markboth{DOC-26 --- Anotações de Oliveira --- O Mecanismo}{}

\noindent \textit{Folhas avulsas encontradas no bolso interno do caderno pessoal do Dr. Heitor Oliveira, amarradas com barbante. Data estimada: junho de 1979 a março de 2026.}

\subsection*{Folha 1 --- 19 de junho de 1979}

\noindent Descobri hoje o que o diário faz com o tempo.

\noindent Escrever é só parte. O diário \textit{endereça}. Quando escrevo no diário de Joaquim, posso escolher \textit{para quando} aquela palavra vai chegar. O controle acontece sem consciência. É como entoar uma nota certa e ver o vidro vibrar. A página vibra. E a palavra some.

\noindent Fiz um teste. Escrevi: ``Maria, se você está lendo isto em 1980, acenda uma vela na janela no dia 15 de agosto.'' Fechei o diário. Guardei debaixo da cama. No dia 15 de agosto de 1979, liguei para a vizinha. Ela disse: ``Sua esposa acendeu uma vela na janela ontem.''

\noindent Funciona.

\noindent O diário não é um registro. É uma \textbf{fenda}. E toda fenda tem uma assinatura.

\subsection*{Folha 3 --- 22 de junho de 1979}

\noindent O padrão apareceu hoje. O diário não aceita qualquer intervalo. Rejeita. A página fica em branco, a tinta escorre, a caneta risca sozinha um risco seco. Funciona apenas no número que ela assinou.

\noindent Sessenta e dois.

\noindent O número de anos que Joaquim passou no Colônia. A entidade gravou o número no diário como uma assinatura. Não sei se é respiração, ciclo ou contagem. Sei que não obedece a mim.

\noindent A página 62 do diário de Joaquim se ilumina quando a palavra vai chegar. É ali que a fenda se abre. Ele escreveu aquela página em 1919, dois anos depois de entrar no Colônia. Já sabia que a fome contava em 62.

\noindent Ele escreveu: ``Sessenta e dois é a idade dela. Ela me disse hoje. Nasceu na seca de 1853, na mesma vala, no mesmo cheiro. Sessenta e dois anos depois, eu nasci. Sessenta e dois anos depois, ela me encontrou.''

\noindent A entidade nasceu em 1853. Nasceu da seca. Da fome. Dos corpos empilhados. Espera 62 anos. Encontra um hospedeiro. Espera mais 62. Outro hospedeiro.

\noindent O ciclo não é humano. É geológico. É o tempo que a terra leva para esquecer.

\subsection*{Folha 15 --- 12 de janeiro de 2026}

\noindent A carta chegou.

\noindent Eu senti. Hoje, 12 de janeiro de 2026, senti a fenda se fechar. A carta foi recebida. Alguém leu. O ciclo encontrou o próximo.

\noindent Olho no espelho e vejo um velho que não sou eu. A pele pendurada. O olho direito completamente dourado. A pupila já não sei se é minha. Ela se alonga quando o quarto escurece. Fica vertical. Fenda de cabra. De bicho que enxerga no escuro.

\noindent Entreguei o diário hoje. Coloquei num pano preto. Deixei na recepção. A mulher da recepção perguntou meu nome. Eu disse: ``Diga que foi o paciente 1.261.''

\noindent Ela não entendeu. Mas a pessoa certa vai entender.

\subsection*{Folha 16 --- 12 de janeiro de 2026, à noite}

\noindent Voltei para casa depois da entrega. Sentei na cama. Olhei minhas mãos. 47 anos. 47 anos vivo contra minha vontade. A carta que escrevi para Maria em 1979 dizia que eu ia proteger os outros. Que eu não ia deixar o diário passar adiante. Acreditei nisso por décadas. Acreditei que estava no controle.

\noindent Não estava. Nunca estive.

\noindent A fome me manteve vivo. Não por bondade. Não por acaso. Ela precisava de um entregador. Alguém que pudesse andar entre os vivos, que tivesse documento, identidade, aparência de gente. Alguém que chegasse numa recepção de arquivo público em 2026 e não levantasse suspeitas. Um velho de 81 anos, de terno puído, com um pacote embaixo do braço. Invisível. Perfeito.

\noindent Ela planejou isso desde o dia em que li o diário, em 1979. Desde o dia em que Joaquim escreveu a página 62, em 1919. Desde o dia em que a vala começou a pulsar, em 1853.

\noindent O ciclo não me perguntou. O ciclo só me usou.

\noindent Eu era só um elo. E elo não escolhe. Elo segura. Elo passa.

\noindent Maria morreu em 2004. Eu não fui ao enterro. Não porque não quis. Porque não \textit{pude}. No dia do enterro, acordei num banco de praça em outra cidade, com as roupas trocadas e 62 dias de memória perdidos. Ela tinha me levado para longe. Ela não ia deixar ninguém interromper a entrega.

\noindent Desculpa, Maria. Desculpa, doutora que vai ler isto. Desculpa, próximo da fila.

\noindent Eu não tive escolha.

\noindent E vocês também não vão ter.

\vspace{0.5em}
\noindent \textit{Nota do arquivista: este maço foi encontrado em fevereiro de 2026, em um envelope pardo lacrado, dentro da mesma caixa que o diário do paciente 1.260. A caligrafia das primeiras folhas é firme, médica. A das últimas é trêmula, com traços que lembram --- mas não idênticos --- à caligrafia de Joaquim. A data da última anotação coincide com a data em que o diário foi devolvido ao Arquivo Público Mineiro pela primeira vez.}

> **↪ Rotas:** **DOC-01** · **LUC-02** · **DOC-27**
